\documentclass[10pt]{article}

% ---------------------------------------------------------------------
% IJCAI-style two-column layout, reproduced from the official documented
% specification (twocolumn, US Letter, 7in x 9in text block, 0.25in
% column separation, no running headers/page numbers, numbered
% citation style) since the exact byte-for-byte ijcai.sty could not be
% retrieved in full in this environment; this is a faithful structural
% match using only packages confirmed available here.
% ---------------------------------------------------------------------
\usepackage[textwidth=7in, textheight=9in, top=1in, left=0.75in,
            columnsep=0.25in, letterpaper]{geometry}
\twocolumn
\usepackage{amsmath}
\usepackage{amssymb}
\usepackage{booktabs}
\usepackage{graphicx}
\usepackage{url}
\usepackage[hidelinks]{hyperref}
\pagestyle{empty}

\setlength{\parindent}{1em}
\setlength{\parskip}{0pt}

% Disable automatic word-hyphenation at line breaks entirely (no word,
% including compound words like "prompt-injection", will ever split
% across a line end), and loosen justification tolerance slightly so
% this doesn't create overfull lines instead.
\hyphenpenalty=10000
\exhyphenpenalty=10000
\tolerance=2000
\emergencystretch=3em

\newcommand{\rs}{\textrm{Rs.\ }}

\begin{document}

\begin{center}
{\large \textbf{Errata: A General-Purpose Evaluation Harness for AI Agents,
Demonstrated on a Bayesian Vendor-Payment-Fraud Triage Agent}}\\[8pt]
{Arif Hussain}\\[10pt]
\end{center}

\noindent\textbf{Abstract.}
    Standard evaluation reports a single number -- ``94\% accurate'' --
    that hides which mistakes a model makes, how expensive each mistake
    actually is, whether the model's stated confidence can be trusted,
    and whether it can be tricked. This paper presents two connected
    contributions. First, a small probabilistic agent that decides
    whether to pay, verify, or escalate a vendor bank-account-change
    request under genuine uncertainty about the sender's identity,
    built with an explicit belief-state table, a hand-derived
    cost-sensitive decision threshold (refined for a three-action
    policy that accounts for imperfect verification), and validated
    through an executed 1{,}000-case Monte Carlo simulation plus five
    stress-test scenarios, including a real reliability audit that
    caught and fixed two genuine implementation bugs. Second, Errata,
    a domain-agnostic evaluation harness intended to audit any such
    agent along five dimensions -- flat accuracy, hierarchical
    tree-distance severity, monetary cost, calibration, and adversarial
    robustness -- validated independently through a dry run on a
    separate demo classifier. We show, concretely, that the
    fraud-triage agent's own confusion-matrix scoring collapses three
    structurally distinct fraud mechanisms into a single class,
    discarding exactly the information Errata's tree-distance scoring
    is designed to recover, and we report this gap, and the
    not-yet-completed integration between the two systems, honestly as
    a limitation and a concrete next step.

\bigskip

\section{Introduction}

Almost every applied machine learning project ends with one number:
``my model is X\% accurate.'' That number is not false, but it is
frequently misleading. Consider a company receiving 100 emails a day,
of which 8 are genuinely dangerous. A filter that says ``safe'' to
every single email, without reading a word, is 92\% accurate --
and catches exactly zero of the 8 dangerous cases. This is the
well-known \emph{accuracy paradox}: when the class of interest is rare,
a trivial policy can score arbitrarily well while providing zero
protection against the thing the system was built to catch.

This problem is not hypothetical. The FBI's Internet Crime Complaint
Center (IC3) reported \$3.05 billion in verified Business Email
Compromise (BEC) losses in 2025 alone, from 24{,}768 complaints, and
60\% of all BEC scams specifically impersonate vendors or
suppliers~\cite{ic3bec}. Named cases include Google and Facebook losing
a combined \$121 million to a fake-vendor invoice
scheme~\cite{proofpointbec}, and the City of Baltimore losing \$1.52
million when an attacker inserted new bank details into an active,
already-trusted email thread~\cite{veridionbec}.

Our central research question is: \emph{how can an AI agent use
probabilistic reasoning to decide what to believe, which evidence is
worth its cost, and when it is rational to act -- and how can a
general-purpose evaluator, independent of the agent's specific domain,
verify that the agent's own self-reported performance can be trusted?}
Narrowed to our concrete problem: \emph{when a supplier requests a
bank-account change, which verification step should the agent perform
next, and at what belief should it stop verifying and either pay or
escalate?}

\section{Related Work}

\textbf{Confusion-matrix evaluation} for flat categories is a mature,
solved problem, implemented correctly by libraries such as
pycm and scikit-learn. We do not reinvent this.

\textbf{Hierarchical, severity-aware scoring} is comparatively rare.
Apple's \emph{Neo} system~\cite{neo2022} treats a confusion matrix as
hierarchical rather than flat, allowing ``mistook malware for
phishing'' to be scored as a smaller error than ``mistook phishing for
safe.'' We initially believed no public implementation of this idea
existed; on independent re-verification we found this claim was
partially wrong -- Neo's \emph{visualization} code is public on
GitHub (\texttt{apple/ml-hierarchical-confusion-matrix}) -- but it
remains a JavaScript/TypeScript rendering library, not a Python
function one can import to compute a tree-distance-aware score inside
a scoring pipeline. That narrower gap is what Errata's own
\texttt{tree\_distance()} implementation fills.

\textbf{Cost-sensitive learning} is well established~\cite{elkan2001},
and scikit-learn now ships
\texttt{TunedThresholdClassifierCV}, which selects a decision threshold
from a cost matrix rather than defaulting to 0.5~\cite{sklearncost}. Our
own threshold derivation (Section~\ref{sec:policy}) is independently
aligned with this official, current tool. What remains largely absent
from the literature is combining cost-sensitivity with hierarchical
tree-distance in a single evaluation report.

\textbf{LLM/agent evaluation harnesses.} AgentSentinel wraps an agent
and checks faithfulness and prompt injection resistance, but its
scoring remains pass/fail, with no notion of cost or hierarchical
severity. DeepEval is a substantially larger, more established
framework (over 100 million daily evaluations reported, adoption
across more than half of the Fortune 500), offering 50+
research-backed metrics, but these remain individual per-criterion
scores rather than a single report combining hierarchy, cost, and
calibration.

\textbf{Adversarial risk.} The OWASP Top 10 for LLM
Applications~\cite{owasp2026} keeps Prompt Injection at \#1 for a
third consecutive year in its 2026 edition (now weighted 75\%
practitioner survey, 25\% analysis of 7{,}714 documented incidents).
The single largest move in the 2026 ranking is \emph{Excessive
Agency}, rising from \#6 to \#3, reflecting that agents are
increasingly trusted with real tools and real financial consequences
-- directly relevant to an agent that can authorize a real payment.

\section{Probabilistic View}

\subsection{Problem Formulation}
The agent observes the email requesting a bank-account change --
sender address, SPF/DKIM result, message content, and whether the
requested account has been used before. It must select
\emph{pay}, \emph{verify by phone}, or \emph{hold and escalate},
because whether the request genuinely originated from the real
supplier or from an attacker is not directly observable.

\subsection{Hidden States and Priors}
Five hidden states are considered, including a residual
``something else'' state, with priors summing to 1.00
(Table~\ref{tab:priors}). All values are stated, hypothetical
assumptions, informed by the BEC literature above rather than measured
company data.

\begin{table}[h]
\centering
\scriptsize
\begin{tabular}{@{}p{2.1in}p{0.5in}@{}}
\toprule
Hidden state & Prior \\
\midrule
Genuine & 0.85 \\
Compromised (mailbox hacked) & 0.06 \\
Impersonation (lookalike domain) & 0.06 \\
Insider / malicious collusion & 0.02 \\
Something else (residual) & 0.01 \\
\bottomrule
\end{tabular}
\caption{Hidden states and priors, summing to 1.00.}
\label{tab:priors}
\end{table}

\subsection{Evidence and Likelihoods}
Three evidence sources were built, each with a full likelihood table
$P(\text{evidence}\mid\text{state})$ for both outcomes:
\textbf{E1} (is the requested account new?), \textbf{E2} (does the
sender domain mismatch the known supplier domain?), and \textbf{E3}
(does the supplier confirm by phone, calling the number already on
file, not the number in the email?). A structurally important
property: E2's likelihood for Compromised is exactly 0 -- if an
attacker has genuinely hijacked the real mailbox, the domain is, by
definition, the real domain, so this check cannot detect mailbox
compromise at all. This matches documented industry findings that BEC
attacks routinely pass SPF/DKIM/DMARC checks, since those protocols
validate sending infrastructure, not human intent.

\subsection{A Full Worked Bayesian Update}
Observing E1 = ``new account,'' the update proceeds by
multiplying each state's prior by its likelihood, summing to obtain
the normalizing constant $P(\text{new}) = 0.1735$, then dividing:

\begin{equation}
P(\text{state}\mid E) = \frac{P(E\mid\text{state})\,P(\text{state})}{\sum_i P(E\mid s_i)\,P(s_i)}
\end{equation}

yielding the posterior (Genuine 0.245, Compromised 0.311,
Impersonation 0.311, Insider 0.104, Something else 0.029), summing to
1.00. A single piece of evidence moved combined fraud probability from
a prior of 15\% to a posterior of \textbf{75.5\%}.

\subsection{Entropy: A Counter-Intuitive Finding}
Shannon entropy $H = -\sum p_i\log_2 p_i$ measures belief spread.
Prior entropy is 0.866 bits; posterior entropy, after the same
evidence above, is \textbf{2.032 bits} -- entropy \emph{increased}.
This runs counter to the common assumption that evidence always
reduces uncertainty: the evidence invalidated the dominant hypothesis
(Genuine) without clearly favoring any single alternative, since
Genuine, Compromised, and Impersonation ended up nearly tied. The
agent moved from confident to confused-but-alert, which is itself
useful information a single accuracy number would never surface.

\section{Information-Theoretic View}

\subsection{Calibration}
An agent is calibrated if, across cases where it states X\% confidence,
roughly X\% actually turn out positive -- a different property from
accuracy. We independently rebuilt the documented model (same priors,
likelihoods, and refined thresholds from Section~\ref{sec:policy}) and
ran 1{,}000 simulated cases; the resulting confusion matrix (TP=148,
FN=1, FP=124, TN=727) reproduced the primary experiment
(Section~\ref{sec:experiment}) exactly, confirming a faithful
reproduction. Bucketing by final stated fraud probability yields an
Expected Calibration Error of \textbf{0.0006} and a Brier Score of
\textbf{0.0165} (0 = perfect; 0.25 = coin-flip). A structurally
interesting finding: all 1{,}000 cases fell into only 2 of 10 possible
confidence buckets (0--10\% and 90--100\%) -- the three-zone threshold
policy pushes every case toward a decisive extreme rather than
lingering in genuine middle-ground uncertainty, a property of the
policy design rather than a flaw in the calibration measurement.

\subsection{Jensen--Shannon Divergence}
Unlike KL divergence, JSD is symmetric and bounded even when one
distribution assigns zero probability to an event the other considers
possible. Computed directly on the documented prior/posterior pair
above, $JSD(\text{prior},\text{posterior}) = 0.288$ bits (JS distance
0.536), confirming quantitatively what the entropy finding suggested:
a single evidence source produced a substantial belief shift. In
production, tracking JSD between a live belief distribution and its
original training-time prior is a practical, bounded distribution-shift
alarm, avoiding KL divergence's failure mode of returning infinity the
first time a genuinely new case type appears.

\subsection{Information Selection}
Table~\ref{tab:infogain} compares expected information gain -- the
prior entropy minus the probability-weighted average posterior entropy
across both outcomes of each evidence source -- against cost.

\begin{table}[h]
\centering
\scriptsize
\begin{tabular}{@{}p{1.15in}p{0.5in}p{0.35in}p{0.75in}@{}}
\toprule
Evidence & Info gain & Cost & Notes \\
\midrule
E1: New account? & 0.347 bits & \rs 0 & Weak on sub-type \\
E2: Domain mismatch? & 0.304 bits & \rs 0 & Blind to Compromised \\
E3: Phone callback & 0.347 bits & \rs 300 & Strongest, not free \\
\bottomrule
\end{tabular}
\caption{Expected information gain vs.\ cost per evidence source.}
\label{tab:infogain}
\end{table}

E1 and E3 are tied as the strongest single checks; E2, despite being
structurally blind to mailbox compromise, remains valuable because it
is highly diagnostic specifically for Impersonation. The policy
implication: run E1 and E2 first, always (free, instant, and
complementary); pay for E3 only if belief remains in the uncertain
zone after the free checks.

\section{Decision Policy}
\label{sec:policy}

\subsection{A Cost-Derived Threshold, and Its Refinement}
With $C_{FP} = \rs 2{,}000$ (wrongly holding a genuine request) and
$C_{FN} = \rs 8{,}00{,}000$ (wrongly paying an actual fraud), the
standard two-action break-even is
$P^* = C_{FP}/(C_{FP}+C_{FN}) = 0.2494\%$ -- markedly below an
intuitive round number, a direct consequence of the 400$\times$ cost
asymmetry.

This formula, however, silently assumes exactly two actions. Our
agent has three (Pay/Verify/Escalate), and Verify is neither free nor
perfect -- E3's own likelihood table states a 5\% chance the callback
is fooled when the true state is Compromised. We therefore derive a
refined three-action threshold accounting for $r$, the
prior-weighted probability that verification itself fails to catch a
fraud that is actually present:

\begin{equation}
P_{low} = \frac{C_V}{C_{FN}(1-r)}, \quad r = 5.33\%
\end{equation}

yielding $P_{low} = 0.0396\%$ -- roughly 6.3$\times$ lower than the
naive two-action figure. Once verification has actually occurred and
its cost is sunk, the remaining choice reverts to a plain two-action
comparison, for which the original, unrefined $P^*=0.2494\%$ is
correct and is reused as $P_{STAR}$ for that later decision point.
$P_{high}=90\%$ remains a stated policy choice, not yet
cost-derived, absent an explicit escalation cost.

\subsection{The Umbrella Problem}
As a sanity check on the general form of this reasoning: with a 35\%
chance of rain, carrying an umbrella costing \rs 20 versus getting wet
costing \rs 200 (a normal walk) yields a 10\% break-even; the same
umbrella against a \rs 1{,}000 wedding outfit yields a 2\% break-even.
The forecast never changes -- only the cost ratio does, exactly
mirroring how the 400$\times$ fraud-to-delay cost ratio pushes the
agent's own threshold far below any intuitive round number.

\section{Experiment}
\label{sec:experiment}

A Monte Carlo simulation (Python, fixed seed=42, isolated
\texttt{random.Random} instance per batch) draws a true hidden state
from the priors, generates E1/E2 from their likelihood tables, runs
the Bayesian update above, applies the refined three-action policy,
generates E3 only when the uncertain zone is entered, and grades the
resulting action against the (only-now-revealed) true state.

\textbf{A genuine reliability audit.} Two real implementation bugs
were found and fixed during development, documented rather than
silently corrected: (1) an early version reused $P_{low}$ for the
post-verification Pay/Escalate decision instead of the correct
$P_{STAR}$, caught via a verbose single-case trace before it could
distort batch results; (2) an early version drew all random numbers
from Python's shared global \texttt{random} module, making exact
reported batch counts silently dependent on unrelated demo code
running first -- fixed by giving every batch its own isolated RNG,
verified by re-running the identical seed after 500 unrelated global
random calls and confirming bit-identical results.

Five stress-test scenarios were executed: (1) base-rate sensitivity
(2\%/15\%/40\% fraud), (2) cost-structure sensitivity, (3)
verification-effectiveness sensitivity (1\%/5\%/20\% miss rate), (4)
threshold-boundary robustness, and (5) adversarial gaming, in which a
simulated attacker deliberately forces E1=known and E2=match on every
fraud case.

\section{Results}

At $n=1{,}000$ (baseline priors, seed=42): confusion matrix TP=148,
FN=1, FP=124, TN=727; Recall = 99.33\%; Precision = 54.41\%; raw
Accuracy = 87.50\% (Table~\ref{tab:results} shows why this number is
misleading); Total Expected Cost = \rs 13{,}28{,}500.

\begin{table}[h]
\centering
\small
\begin{tabular}{lc}
\toprule
Metric & Value \\
\midrule
Recall & 99.33\% \\
Precision & 54.41\% \\
Accuracy (misleading) & 87.50\% \\
Verification rate & 93.5\% \\
Total Expected Cost & \rs 13,28,500 \\
\bottomrule
\end{tabular}
\caption{Baseline simulation results ($n=1{,}000$).}
\label{tab:results}
\end{table}

Stress scenarios 1 and 3 independently expose the same underlying
gap: $P_{STAR}$ is a \emph{fixed} cost-ratio number that does not
adapt to base rate or verification quality. At a 40\% fraud base rate,
or a 20\% verification miss rate, even a genuine case with the
best-possible evidence combination computes a final fraud probability
(0.80\% and 0.336\% respectively) that still exceeds the fixed
$P_{STAR}=0.2494\%$, causing near-universal false-positive escalation
in those regimes -- a genuine, simulation-confirmed limitation of a
fixed threshold, hand-verified independently.

\textbf{Comparison against Errata's own dry-run validation.} On a
separate, independently validated demo (22 hand-mocked cases on a
safe/suspicious/dangerous tree), Errata's own scoring reported 77.3\%
flat accuracy alongside a tree-distance of 0.545, a cost of
\rs 90{,}91{,}136 per 100 cases, a calibration error of 0.256, and a
50\% adversarial pass rate -- demonstrating that the same
underlying methodology (flat score, tree distance, cost, calibration,
adversarial testing) produces a materially richer report than a single
accuracy figure, independent of which concrete domain it is applied
to.

\section{Failure Analysis}

Beyond the two reliability bugs above (classified as \emph{wrong
threshold reuse} and \emph{implementation fragility}, both caught and
corrected before they could distort reported results), a specific,
repeatable pattern emerged in the baseline batch's 124 false
positives: cases where E1/E2/E3 all point toward Genuine still leave
several percent of residual belief above the very low $P_{low}$, so
the policy consistently routes them to Verify rather than Pay. This is
classified as \emph{wrong threshold} rather than \emph{bad evidence} or
\emph{bad prior} -- it is a direct, correctly-derived consequence of
the 400$\times$ cost asymmetry, not a bug, and is reported as the
central trade-off finding of the experiment rather than something to
be silently tuned away.

A further gap, found by direct comparison against Errata's design
(rather than by re-running any code): the base-model's own confusion
matrix collapses Compromised, Impersonation, and Insider into a single
``Fraud'' class. Applying Errata's tree-distance scoring to the same
simulation output -- available as a saved, per-case
\texttt{results.csv}, independently reconfirmed to reproduce the
baseline confusion matrix exactly -- would reveal whether the
agent's false positives and false negatives cluster around specific
sub-type confusions, information currently invisible in the flat
matrix.

\section{Human Discussions}

Public discussion (Reddit, X) about this project's problem framing and
prior-art gap is tracked separately in \texttt{discussion-record.md},
as a continuing, live activity rather than a fixed set of quotes to be
reproduced here; at the time of writing this section is a
work-in-progress rather than complete, and is reported as such rather
than populated with any invented exchange.

\section{Limitations, Ethics, and Human Control}

All priors and costs are stated, hypothetical assumptions, not
measured from real company data. The phone-callback likelihood table
assumes only the email channel, never the phone channel, is
compromised; a sufficiently sophisticated attacker who compromises
both would defeat this check entirely. $P_{high}=90\%$ is a policy
choice, not yet cost-derived. $P_{STAR}$ is a fixed cost-ratio number,
simulation-confirmed not to adapt to base rate or verification
quality. Most importantly: \textbf{Errata and the fraud-triage agent
were developed and independently validated in parallel, but have not
yet been integrated into a single pipeline} -- the agent's own
confusion-matrix code, not Errata's tree-distance/cost/calibration
scoring, produced the results in Section~\ref{sec:experiment}. This
integration, using the already-exported \texttt{results.csv}, is
reported honestly as immediate future work rather than implied as
complete. On ethics: an agent authorizing real financial payments
requires human escalation above a stated risk threshold regardless of
confidence, which this design implements; no personally identifying
information beyond what a real accounts-payable workflow already
handles is modeled or stored.

\section{New Questions}

(1) Would Errata's tree-distance scoring, applied to the exported
\texttt{results.csv}, reveal structured sub-type confusion currently
invisible to the base-model's flat confusion matrix? (2) Can
$P_{high}$ be derived from an explicit escalation cost the same way
$P_{low}$ was refined for verification cost? (3) Does a fixed
$P_{STAR}$ generalize acceptably across the base-rate range actually
observed in a real deployment, or does the base-model need the same
adaptive-threshold treatment $P_{low}$ already received?

\section{AI-Use Statement}

An AI assistant (Claude, via Microsoft Copilot) was used throughout
this project to: explain probability, information-theoretic, and
calibration concepts in plain language on request; check arithmetic
and re-derive every reported number independently before it was
trusted (two genuine numeric errors were found this way and corrected
in the decision record rather than silently fixed); debug the
simulation script; and assist with LaTeX formatting for this document.
The hidden states, priors, costs, evidence likelihood tables, decision
thresholds, and the interpretation of every result are the author's
own, arrived at through guided questioning rather than being generated
wholesale. All experiments reported were actually run by the author
(with the assistant verifying reproducibility), using a fixed,
reported random seed.

\section{Conclusion}

We presented a probabilistic vendor-payment-fraud triage agent,
built and validated with real Bayesian reasoning, a cost-derived
and subsequently refined decision threshold, and an executed,
reliability-audited simulation; and Errata, a general-purpose
evaluation harness intended to audit any such agent along five
dimensions beyond flat accuracy. We showed concretely where the two
currently diverge -- the agent's own scoring collapses information
Errata is designed to preserve -- and reported the remaining
integration work honestly rather than overstating current completeness.

\bibliographystyle{plain}
\begin{thebibliography}{9}

\bibitem{ic3bec} Internet Crime Complaint Center (IC3), ``Business
Email Compromise: The \$55 Billion Scam,'' FBI Public Service
Announcement, 2024/2025 updates.

\bibitem{proofpointbec} Proofpoint, ``10 Real-World Business Email
Compromise (BEC) Scam Examples,'' 2023.

\bibitem{veridionbec} Veridion, ``Vendor Fraud Schemes: 5 Real-Life
Examples to Learn From,'' 2026.

\bibitem{neo2022} J.\ G\"ortler, F.\ Hohman, D.\ Moritz, K.\
Wongsuphasawat, D.\ Ren, R.\ Nair, M.\ Kirchner, K.\ Patel, ``Neo:
Generalizing Confusion Matrix Visualization to Hierarchical and
Multi-Output Labels,'' ACM CHI, 2022.

\bibitem{elkan2001} C.\ Elkan, ``The Foundations of Cost-Sensitive
Learning,'' IJCAI, 2001.

\bibitem{sklearncost} scikit-learn developers, ``Post-tuning the
Decision Threshold for Cost-Sensitive Learning,''
\texttt{TunedThresholdClassifierCV} documentation, 2025.

\bibitem{owasp2026} OWASP GenAI Security Project, ``OWASP Top 10 for
LLM Applications,'' 2026 edition.

\end{thebibliography}

\onecolumn
\appendix
\section*{Appendix A: Required Core Questions}
{\small
\begin{enumerate}
\item \textbf{Hidden state:} an unobservable fact -- here, the true
identity/intent behind a bank-change request. Ours: Genuine,
Compromised, Impersonation, Insider, Something else.
\item \textbf{Why not predict directly:} a single point prediction
discards the agent's own uncertainty, which is precisely the signal
needed to decide when to escalate to a human.
\item \textbf{Belief distribution:} a probability over all hidden
states; must sum to one because the states are exhaustive and mutually
exclusive.
\item \textbf{Prior vs.\ posterior:} prior is belief before evidence
(0.85 / 0.06 / 0.06 / 0.02 / 0.01); posterior is belief after
(0.245 / 0.311 / 0.311 / 0.104 / 0.029, given E1=new).
\item \textbf{Likelihood vs.\ probability:} likelihood is
$P(\text{evidence}\mid\text{state})$, read state-first; it is not the
posterior, which is evidence-first.
\item \textbf{Agent uncertainty:} the agent holds several explanations
simultaneously rather than collapsing to one guess.
\item \textbf{Entropy:} $H=-\sum p\log_2 p$; ours rose from 0.866 to
2.032 bits after one evidence source -- a genuine, documented
exception to ``evidence always reduces uncertainty.''
\item \textbf{A bit:} one unit of entropy, roughly one good yes/no
question's worth of uncertainty.
\item \textbf{Information gain:} $H(\text{prior})-H(S\mid E)$; E1 and
E3 both scored 0.347 bits, E2 scored 0.304.
\item \textbf{Conditional entropy:} the probability-weighted average
of posterior entropy across all possible evidence outcomes.
\item \textbf{Mutual information vs.\ correlation:} the same quantity
as information gain; unlike correlation, it captures non-linear
dependence (e.g., a variable and its square have zero correlation but
maximal mutual information).
\item \textbf{KL divergence:} not a true distance -- asymmetric, and
infinite if one distribution assigns zero probability to an event the
other considers possible.
\item \textbf{Asymmetry example:} comparing a fair coin against a
model that says ``always heads'' gives infinite divergence one
direction, finite the other.
\item \textbf{Jensen--Shannon divergence:} the symmetric, bounded fix,
comparing both distributions to their average; computed here as 0.288
bits between our prior and posterior.
\item \textbf{Calibration:} whether stated confidence matches observed
frequency; ours scored ECE=0.0006, Brier=0.0165 -- excellent, though
concentrated in only two of ten confidence buckets.
\item \textbf{Expected cost / threshold:} the action minimizing
expected cost, not the most probable state; ours derived $P^*=$
0.2494\%, later refined to $P_{low}=0.0396\%$ for the three-action
case.
\item \textbf{Value of information:} evidence is only worth its cost
if some outcome would change the action; if every outcome leads to the
same decision, its value is zero regardless of its bit count.
\item \textbf{Highly informative but useless:} E2 cannot distinguish
Genuine from Compromised at all (both 0.00 likelihood for mismatch),
yet still carries real information for Impersonation specifically --
demonstrating information value is state-specific, not monolithic.
\item \textbf{Distribution shift:} the world changing since priors
were set; a rising JSD between a live belief distribution and the
original prior is a practical, bounded detection signal.
\item \textbf{Stop rule:} act once no remaining check could change the
decision, the belief has left the uncertain zone, or the next check's
cost exceeds its expected benefit.
\end{enumerate}
}

\end{document}
